\setcounter{theorem}{40}
\appendix
\section{Karger's theorem}
\label{app:karger}
In this appendix we prove a slight generalization of Karger's theorem \cite[Theorem 4.1]{Karger00} which is needed for our time-efficient algorithm.  
We begin by introducing some needed tools.

\subsection{Tools}
Matula \cite{Mat93} gave an $O(m/\varepsilon)$ time deterministic algorithm to compute a $(2+\varepsilon)$-approximation to the edge connectivity of a simple graph (or multigraph).  
The algorithm can also be adapted to give a constant factor approximation to the weight of a minimum cut in an integer-weighted graph in time $O(m \log^2(n))$, see Appendix~A of \cite{GMW20}.
\begin{lemma}[Matula's approximation algorithm \cite{Mat93, GMW20}]
\label{lem:matula}
Let $G = (V,w)$ be an integer-weighted graph with $m$ edges and $n$ vertices.  There is a constant $c$ and a deterministic algorithm that in time $O(m \log^2(n))$
outputs a value $\tilde \lambda$ such that $\tilde \lambda/c \le \lambda(G) \le  \lambda$.
\end{lemma}

To efficiently construct a tree-packing we will also need to use random sampling.  The following lemma is the heart of Karger's skeleton construction \cite{Karger99}.  We recommend 
the presentation in \cite[Lemma 14]{BLS20}.
\begin{lemma}[\cite{Karger99}]
\label{lem:downsample}
Let $G$ be an unweighted multigraph with $m$ edges.  For an integer $d \ge 2$ and real numbers $\varepsilon, \gamma$ with $\varepsilon \le 1/3$, 
let $p = 3d(\ln n)/(\varepsilon \lambda(G))$.  In time $O(pm \log(n))$  we can randomly sample $\ceil{pm}$ edges of $G$.  With probability $1-1/n^d$ the resulting graph $H$ has the properties that
\begin{enumerate}
  \item The minimum cut of $H$ is within a $(1+\epsilon)$ factor of $p\lambda(G) = 3d\ln(n)/\varepsilon^2$.
  \item For every $X \subseteq V$ we have $(1-\varepsilon) w(\Delta_G(X)) \le w(\Delta_H(X)) \le (1+\varepsilon) w(\Delta_G(X))$.
\end{enumerate}
\end{lemma}

Another very useful tool we use is the Nagamochi-Ibaraki construction which shows that for an integer-weighted graph $G$ with $m$ edges, in time $O(m \log(n))$ one can construct 
a graph $G'$ whose total edge weight is $nc$ and which preserves all cuts of $G$ of weight at most $c$.  

\begin{lemma}[\cite{NI92}]
Let $G =(V,w)$ be an $n$-vertex integer-weighted graph with $m$ edges.  For any positive integer $c$ there is a deterministic algorithm that in time $O(m \log n)$ produces an integer-weighted graph $G' = (V,w')$ with 
total edge weight $O(cn)$ such that for all $X \subseteq V$ with $\Delta_G(X) \le c$ it holds that $w(e) = w'(e)$ for all $e \in \Delta_G(X)$.  Thus in particular $\Delta_G(X) = \Delta_{G'}(X)$ and $w(\Delta_G(X)) = w'(\Delta_{G'}(X))$ for 
all $X$ with $\Delta_G(X) \le c$.
\end{lemma}

We combine the tools of Matula's approximation algorithm, random sampling, and the sparse certificate of Nagamochi-Ibaraki into the following lemma.
\begin{lemma}
\label{lem:Hprop}
Let $G = (V,w)$ be an integer-weighted graph and let $0 < \delta < 1$ be a parameter.  There is an $O(m \log^2(n) + n\log(n))$ time randomized algorithm to create a weighted graph $H = (V, w_H)$ such that 
\begin{enumerate}
  \item $H$ has $O(n \log(n)/\varepsilon^2)$ edges.
  \item The minimum cut of $H$ has value $\lambda(H) = O(\log n)$.
  \item If $X \subseteq V$ is such that $w(\Delta(X)) \le (1+\delta) \lambda(G)$ then $w_H(\Delta(X)) \le (1+3\delta) \lambda(H)$.
\end{enumerate}
\end{lemma}

\begin{proof}
First, by \cref{lem:matula}, in time $O(m \log^2(n))$ we can find a constant factor approximation $\tilde \lambda$ satisfying $\tilde \lambda/c \le \lambda(G) \le \tilde \lambda$ .  Next we apply the Nagamochi-Ibaraki 
algorithm to $G$ with threshold $t = (1+\delta)\tilde \lambda$.  In $O(m \log(n))$ time this produces an integer-weighted graph $G_2 = (V,w')$ with total edge weight $O(tn)$ such that for every $X \subseteq V$ with 
$w(\Delta_G(X)) \le (1+\delta) \tilde \lambda$ it holds that $w(\Delta_G(X)) = w'(\Delta_G(X))$.  

We now view $G_2$ as an unweighted multigraph with $O(tn)$ edges and apply \cref{lem:downsample}.
Let $p = \ln(n)/\tilde \lambda$.  We randomly choose $\ceil{p E(G_2)} =  O(n \ln(n))$ edges of $G_2$ and let the resulting graph be $H$.  This can be done in time $O(n \log(n))$.  
By \cref{lem:downsample} the graph has the stated properties.  The total running time is $O(m\log^2(n)+ n \log(n))$.
\end{proof}

\subsection{Tree packing}
With these preliminaries in place we now turn to actually constructing a tree packing.  We first need the definition, and a lemma of Karger.
\begin{definition}[Weighted tree packing]
Let $G = (V,w)$ be an integer-weighted graph.  A \emph{weighted tree packing} is a set of spanning trees of $G$, each with an assigned weight, such that 
the total weight of trees containing any edge $e \in E(G)$ is at most $w(e)$.  The \emph{value} of the packing is the total weight of trees in it.  
\end{definition}

\begin{lemma}[{\cite[Lemma 2.3]{Karger00}}]
\label{lem:frac_trees}
Given a weighted tree packing of value $\beta c$ and a cut of value $\alpha c$, at least a $(3-\alpha/\beta)/2$ fraction of the trees by weight 2-constrain the cut.
\end{lemma}

Gabow gives an algorithm to construct a near optimal tree packing in an unweighted multigraph.  The following is an easy adaptation to an integer-weighted graph.
\begin{lemma}[\cite{Gabow95}]
\label{lem:tree_packing}
Let $G = (V,w)$ be an integer-weighted graph with $n$ vertices and $m$ edges.  There is a deterministic algorithm that finds an integer-weighted tree packing of $G$ of value at least 
$\lambda(G)/2$ in time $O(m (\lambda(G)^2 \log(n) + \log^2(n)))$.  
\end{lemma}

\begin{proof}
For a multigraph $H$ with $n$ vertices and $m'$ edges, Gabow \cite{Gabow95} gives a deterministic algorithm that finds a tree packing of weight $\lambda(H)/2$ in time 
$m' \lambda(H) \log(n)$.  The only difference with our case is that $G$ is an integer-weighted graph instead of a multigraph.  We can of course view $G$ as a multigraph but it 
becomes too expensive to run Gabow's algorithm if this significantly blows up the number of edges.  

Thus we first use \cref{lem:matula} to compute $\tilde \lambda$ such that $\tilde \lambda/c \le \lambda(G) \le \tilde \lambda$ in time $O(m \log^2(n))$.  Then we make a pass through the 
edges of $G$ and form a graph $G'$ where any edge of weight larger than $\tilde \lambda$ in $G$ is thresholded down to $\tilde \lambda$.  Thus when viewed as a multigraph $G'$ will only have $O(m \lambda(G))$ edges.  Any tree packing of $G$ is also a tree packing of $G'$ as the value of any tree packing is at most $\lambda(G) \le \tilde \lambda$.  We can then apply Gabow's algorithm to $G'$ to obtain the theorem.
\end{proof}

We are finally ready to prove the slight generalization of Karger's theorem that we require.
\kargerthm*

\begin{proof}
In $O(m)$ time we can find the minimum weight $\alpha$ of an edge of $G$.  Multiplying all edge weights by $1/\alpha$ we obtain a graph where all edge weights are at least $1$ and that has the same 
set of $(1+1/16)$-near minimum cuts as $G$.  Thus without loss of generality now assume that $G$ has all edge weights at least $1$.  

In $O(m)$ time we create the integer-weighted graph $G' = (V,w')$ where $w'(e) = \lfloor 100w(e) \rceil$.  Note that as we assume that every edge of $G$ has weight at least $1$, for any 
$X \subseteq V$ we have
\begin{equation}
\label{eq:int_approx}
0.995 w(\Delta_G(X)) \le \frac{ w(\Delta_{G'}(X))}{100} \le 1.005 w(\Delta_G(X)) \enspace .
\end{equation}
Thus if $\Delta_G(X)$ is a $(1+ \varepsilon)$-near minimum cut of $G$ then $\Delta_{G'}(X)$ is a $(1+\varepsilon)(1.005)^2$-near minimum cut of $G'$.  With $\varepsilon = 1/16$ 
it follows that $\Delta_{G'}(X)$ is a $1+1/12$-near minimum cut of $G'$.

Next we apply \cref{lem:Hprop} to $G'$ to in time $O((m + n) \log^2(n))$ create a graph $H$ with the properties specified there.  We then use \cref{lem:tree_packing} to find a tree packing 
of weight at least $\lambda(H)/2$ and which contains $O(\log(n))$ trees since $\lambda(H) = O(\log(n))$.  Now let $\Delta_G(X)$ be a $(1+1/16)$-near minimum cut of $G$. 
Then $\Delta_{G'}(X)$ is a $(1+1/12)$-near mincut of $G'$ and by \cref{lem:Hprop}, $\Delta_H(X)$ is a $1+1/4$-near mincut of $H$.  
Therefore by \cref{lem:frac_trees} at least $1/4$ of the trees in the packing will $2$-respect $\Delta_H(X)$.  
These trees must also $2$-respect $\Delta_G(X)$ since it has the same shore $X$.
\end{proof}

